%%%%%%%%%%%%%%%%%%%%%%%%%%%%%%%%%%%%%%%%%%%%%%%%%%%%%%%%%%%%%%%%%%%%%%
% Lớp 8 | Chương I. Đa thức | Phép chia đa thức cho đơn thức
% Dạng 8D1?5-3: Xét tính chia hết, tìm điều kiện để phép chia hết
%%%%%%%%%%%%%%%%%%%%%%%%%%%%%%%%%%%%%%%%%%%%%%%%%%%%%%%%%%%%%%%%%%%%%%


% nguồn: De_AI_Toan8_cuoi_hoc_ki_i_1 (nhập 21/11/2026)
\begin{bt}%[8D1V5-3]%[ma:59e4aa75]
	Tìm số tự nhiên $n$ để mỗi phép chia sau là phép chia hết:
a) $x^4y^3 : (x^ny^2)$;
b) $6x^ny^5 : (2x^3y^n)$.
	\loigiai{
		a) Để $x^4y^3$ chia hết cho $x^ny^2$ thì số mũ của mỗi biến trong đơn thức bị chia không nhỏ hơn số mũ tương ứng trong đơn thức chia: \\
$n \le 4$ và $3 \ge 2$ (luôn đúng). Vì $n \in \mathbb{N}$ nên $n \in \{0; 1; 2; 3; 4\}$. \\
b) Để $6x^ny^5$ chia hết cho $2x^3y^n$ thì $n \ge 3$ và $n \le 5$, tức $3 \le n \le 5$. \\
Vì $n \in \mathbb{N}$ nên $n \in \{3; 4; 5\}$.
	}
\end{bt}

% nguồn: De_AI_Toan8_cuoi_hoc_ki_i_1 (nhập 21/11/2026)
\begin{bt}%[8D1V5-3]%[ma:fd1ee6e5]
	Tìm số tự nhiên $n$ để đơn thức $5x^{n+1}y^6$ chia hết cho đơn thức $3x^4y^{2n}$, rồi thực hiện phép chia với $n$ vừa tìm được.
	\loigiai{
		Để $5x^{n+1}y^6$ chia hết cho $3x^4y^{2n}$ thì $n + 1 \ge 4$ và $6 \ge 2n$. \\
Suy ra $n \ge 3$ và $n \le 3$, do đó $n = 3$ (thỏa mãn $n \in \mathbb{N}$). \\
Khi $n = 3$, phép chia trở thành $5x^4y^6 : (3x^4y^6) = \dfrac{5}{3}$.
	}
\end{bt}

% nguồn: De_AI_Toan8_cuoi_hoc_ki_i_1 (nhập 21/11/2026)
\begin{bt}%[8D1V5-3]%[ma:5bdcabf2]
	Tìm số tự nhiên $n$ để đa thức $A = 5x^3y^4 - 3x^4y^3 + x^5y^2$ chia hết cho đơn thức $B = 2x^ny^n$.
	\loigiai{
		Để đa thức $A$ chia hết cho đơn thức $B$ thì mọi hạng tử của $A$ đều phải chia hết cho $B$: \\
- Xét biến $x$: $n \le 3$, $n \le 4$, $n \le 5 \Rightarrow n \le 3$. \\
- Xét biến $y$: $n \le 4$, $n \le 3$, $n \le 2 \Rightarrow n \le 2$. \\
Kết hợp hai điều kiện trên ta được $n \le 2$. Vì $n \in \mathbb{N}$ nên $n \in \{0; 1; 2\}$.
	}
\end{bt}

% nguồn: De_AI_Toan8_cuoi_hoc_ki_i_1 (nhập 21/11/2026)
\begin{bt}%[8D1V5-3]%[ma:03bfd83d]
	Tìm số tự nhiên $n$ để đa thức $A = 7x^{n-1}y^5 - 5x^3y^4$ chia hết cho đơn thức $B = 5x^2y^n$.
	\loigiai{
		Điều kiện để đa thức $A$ có nghĩa là $n - 1 \ge 0 \Rightarrow n \ge 1$. \\
Để $A$ chia hết cho $B$, từng hạng tử của $A$ phải chia hết cho $B$: \\
- Hạng tử $7x^{n-1}y^5$ chia hết cho $5x^2y^n$ khi $n - 1 \ge 2$ và $5 \ge n \Rightarrow 3 \le n \le 5$. \\
- Hạng tử $-5x^3y^4$ chia hết cho $5x^2y^n$ khi $3 \ge 2$ (đúng) và $4 \ge n \Rightarrow n \le 4$. \\
Kết hợp lại ta được $3 \le n \le 4$. Vì $n \in \mathbb{N}$ nên $n \in \{3; 4\}$.
	}
\end{bt}

% nguồn: De_AI_Toan8_cuoi_hoc_ki_i_1 (nhập 21/11/2026)
\begin{bt}%[8D1V5-3]%[ma:57a0cd05]
	Tìm tất cả các số tự nhiên $n$ để đa thức $P = 12x^6y^4 + 9x^5y^3 - 15x^4y^n$ chia hết cho đơn thức $Q = 3x^ny^2$.
	\loigiai{
		Để $P$ chia hết cho $Q$, mỗi hạng tử của $P$ phải chia hết cho $Q$: \\
- Xét biến $x$: $n \le 6$, $n \le 5$, $n \le 4 \Rightarrow n \le 4$. \\
- Xét biến $y$: $4 \ge 2$ (đúng), $3 \ge 2$ (đúng) và $n \ge 2$. \\
Kết hợp hai điều kiện ta được $2 \le n \le 4$. Vì $n \in \mathbb{N}$ nên $n \in \{2; 3; 4\}$.
	}
\end{bt}

% nguồn: De_AI_Toan8_cuoi_hoc_ki_i_1 (nhập 21/11/2026)
\begin{bt}%[8D1V5-3]%[ma:34b2bc2d]
	Tìm tất cả các số tự nhiên $n$ để đa thức $A = 4x^{n+2}y^3 - 6x^5y^{n+1}$ chia hết cho đơn thức $B = 2x^3y^n$.
	\loigiai{
		Để $A$ chia hết cho $B$, từng hạng tử của $A$ phải chia hết cho $B$: \\
- Hạng tử $4x^{n+2}y^3$ chia hết cho $2x^3y^n$ khi $n + 2 \ge 3$ và $3 \ge n \Rightarrow 1 \le n \le 3$. \\
- Hạng tử $-6x^5y^{n+1}$ chia hết cho $2x^3y^n$ khi $5 \ge 3$ và $n + 1 \ge n$ (luôn đúng với mọi $n \in \mathbb{N}$). \\
Vậy $n \in \{1; 2; 3\}$.
	}
\end{bt}

% nguồn: De_AI_Toan8_cuoi_hoc_ki_i_1 (nhập 21/11/2026)
\begin{bt}%[8D1V5-3]%[ma:dbec04d4]
	Tìm số nguyên dương $n$ để đa thức $M = 2x^5y^3 - 4x^3y^4 + 6x^2y^2$ chia hết cho đơn thức $N = x^ny^{n-1}$.
	\loigiai{
		Với $n$ nguyên dương ($n \ge 1$), để $M$ chia hết cho $N$ thì: \\
- Xét biến $x$: $n \le \min(5; 3; 2) = 2$. \\
- Xét biến $y$: $n - 1 \le \min(3; 4; 2) = 2 \Rightarrow n \le 3$. \\
Kết hợp với $n \ge 1$ và $n \in \mathbb{Z}$, ta được $n \in \{1; 2\}$.
	}
\end{bt}

% nguồn: De_AI_Toan8_cuoi_hoc_ki_i_1 (nhập 21/11/2026)
\begin{bt}%[8D1V5-3]%[ma:1261b22e]
	Tìm tất cả các số tự nhiên $n$ để phép chia $(3x^{n+1}y^5 - 5x^4y^{n+1}) : (x^3y^n)$ là phép chia hết.
	\loigiai{
		Để phép chia là phép chia hết, từng hạng tử của đa thức bị chia phải chia hết cho $x^3y^n$: \\
- Hạng tử $3x^{n+1}y^5$ chia hết cho $x^3y^n$ khi $n + 1 \ge 3$ và $5 \ge n \Rightarrow 2 \le n \le 5$. \\
- Hạng tử $-5x^4y^{n+1}$ chia hết cho $x^3y^n$ khi $4 \ge 3$ và $n + 1 \ge n$ (luôn đúng). \\
Vì $n \in \mathbb{N}$ nên $n \in \{2; 3; 4; 5\}$.
	}
\end{bt}

% nguồn: De_AI_Toan8_cuoi_hoc_ki_i_1 (nhập 21/11/2026)
\begin{bt}%[8D1V5-3]%[ma:dd8283a7]
	Tìm số nguyên dương $n$ để đa thức $A = 10x^4y^3z^2 - 5x^3y^4z^3 + 15x^2y^2z^4$ chia hết cho đơn thức $B = 5x^ny^nz^n$. Với $n$ lớn nhất tìm được, hãy thực hiện phép chia $A : B$.
	\loigiai{
		Để $A$ chia hết cho $B$ thì với mọi hạng tử: \\
- Xét biến $x$: $n \le \min(4; 3; 2) = 2$. \\
- Xét biến $y$: $n \le \min(3; 4; 2) = 2$. \\
- Xét biến $z$: $n \le \min(2; 3; 4) = 2$. \\
Vì $n$ nguyên dương nên $n \in \{1; 2\}$. Giá trị lớn nhất là $n = 2$. \\
Khi $n = 2$, $A : B = (10x^4y^3z^2 - 5x^3y^4z^3 + 15x^2y^2z^4) : (5x^2y^2z^2) = 2x^2y - xy^2z + 3z^2$.
	}
\end{bt}

% nguồn: De_AI_Toan8_cuoi_hoc_ki_i_1 (nhập 21/11/2026)
\begin{bt}%[8D1V5-3]%[ma:2b6d0829]
	Tìm các cặp số tự nhiên $(m; n)$ thỏa mãn $m + n = 7$ sao cho đa thức $A = 3x^my^4 - 7x^5y^n$ chia hết cho đơn thức $B = 2x^3y^2$.
	\loigiai{
		Để $A$ chia hết cho $B$, từng hạng tử của $A$ phải chia hết cho $2x^3y^2$: \\
- Hạng tử $3x^my^4$ chia hết cho $2x^3y^2$ khi $m \ge 3$ (vì $4 \ge 2$ đã đúng). \\
- Hạng tử $-7x^5y^n$ chia hết cho $2x^3y^2$ khi $n \ge 2$ (vì $5 \ge 3$ đã đúng). \\
Vì $m, n \in \mathbb{N}$, $m + n = 7$, $m \ge 3$ và $n \ge 2$, ta có các cặp $(m; n)$ thỏa mãn là $(3; 4)$, $(4; 3)$, $(5; 2)$.
	}
\end{bt}

% nguồn: De_AI_Toan8_cuoi_hoc_ki_i_1 (nhập 21/11/2026)
\begin{bt}%[8D1V5-3]%[ma:b67d0e49]
	Tìm số tự nhiên $n$ nhỏ nhất để đa thức $A = 8x^{2n}y^5 - 12x^{n+3}y^{n+1}$ chia hết cho đơn thức $B = 4x^6y^4$. Thực hiện phép chia $A : B$ với giá trị $n$ vừa tìm được.
	\loigiai{
		Để $A$ chia hết cho $B$, từng hạng tử của $A$ phải chia hết cho $4x^6y^4$: \\
- Hạng tử $8x^{2n}y^5$ chia hết cho $4x^6y^4$ khi $2n \ge 6 \Rightarrow n \ge 3$ (vì $5 \ge 4$ đã đúng). \\
- Hạng tử $-12x^{n+3}y^{n+1}$ chia hết cho $4x^6y^4$ khi $n + 3 \ge 6$ và $n + 1 \ge 4 \Rightarrow n \ge 3$. \\
Suy ra $n \ge 3$, nên số tự nhiên $n$ nhỏ nhất là $n = 3$. \\
Với $n = 3$, ta có $A : B = (8x^6y^5 - 12x^6y^4) : (4x^6y^4) = 2y - 3$.
	}
\end{bt}
